\chapter{Méthodes d'estimation}

\section{Exemple introductif : Exemple du sondage}

Un projet de loi est soumis à referendum. La proportion $\theta$ d'individu qui vont voter "oui" est un réel inconnu de $[0,1]$. Pour savoir si le oui va l'emporter, on réalise un sondage et on interroge $n$ individus dans la population française. On tire ces individus au hasard et avec remplacement.

On note $X_i = 1$ si la réponse du $i$-éme individu indépendantes et identiquement distribuées, de loi de Bernoulli $\mathcal{B}(\theta)$. Au vu des réalisations des variables aléatoires $X_i$, on cherche à déterminer la valeur de $\theta$ et savoir si le projet de loi va être adopté : on souhaite \textit{estimer} $\theta$ et \textit{tester} si $\theta$ est supérieur ou non à $\frac{1}{2}$. Il est naturel de considérer comme estimateur de $\theta$ la moyenne empirique:

$$
\hat{\theta} = \frac{1}{n} \sum_{i=1}^n X_i
$$  

L'inégalité de Tchebychev nous dit que, pour tout $\varepsilon >0$,

$$
\mathbb{P}_\theta \left( |\hat{\theta} - \theta | \geq \varepsilon \right) \leq \frac{\mathbb{V}[X_i]}{n \varepsilon ^2} \leq \frac{1}{4 n \varepsilon^2}
$$

Noter que la loi de $\theta$ dépend du $\theta$ inconnu, d'où l'indice dans la probabilité.

Puisque la variance d'une loi de Bernoulli est $\theta (1-\theta)$ et est majorée uniformément par $\frac{1}{4}$ le choix de $\varepsilon = \frac{1}{2 \sqrt{n \alpha}}$ conduit à :

$$
\mathbb{P}_\theta \left( |\hat{\theta} - \theta | \leq \frac{1}{2 \sqrt{n \alpha}} \right) \geq 1 - \alpha
$$

On dit que $[\hat{\theta} - \frac{1}{2 \sqrt{n \alpha}} ; \hat{\theta}  + \frac{1}{2 \sqrt{n \alpha}}]$ est un intervelle de condiance pour $\theta$ de niveeau de confiance $1 - \alpha$.

Pour $\alpha = 5 \%$ et $n=1000$, la demi-longuer de l'intervalle est environ égale à $0.0707$.

Ce résultat montre un biais fréquent dans l'interprétation des sondages par les médias. Par exemple, si un journal annonce une variation de $2 \%$ des intentions de vote au cours de $6$ mois sur la base d'un échantillon de $1000$ personnes, ce changement n'a aucune valeur. Comme on a montré, ce mouvement apparent n'est que du bruit statistique.

\newpage

\section{Modèles statistiques}

\begin{definition}
On appelle \textbf{observation} la variable aléatoire $X = ( X_1 , \cdots , X_n)$

On dit que $( X_1 , \cdots , X_n)$  est un $n$-\textbf{échantillon} de loi $P$ si les variables aléatoires $X_1 , \cdots , X_n$ sont indépendantes identiquement distrubuées (i.i.d.) de loi $P$.
\end{definition}

\smallskip

\begin{mdframed}
\begin{definition}[Modèle statistique]
Un modèle statistique est un triplet :
\[
(E, \mathcal{E}, (P_\theta)_{\theta \in \Theta})
\]
où :
\begin{itemize}
  \item $(E, \mathcal{E})$ est un espace mesurable probabilisable (dans ce cours, on retiendra surtout que $X \in E$) ;
  \item $(P_\theta)_{\theta \in \Theta}$ est une famille de lois de probabilité sur $E$ ;
  \item $\Theta$ est l'ensemble des paramètres.
\end{itemize}

Si $\Theta \subset \mathbb{R}^d$ avec $d \geq 1$, le modèle est dit \textbf{paramétrique}.  
Sinon, il est dit \textbf{non paramétrique}.
\end{definition}
\end{mdframed}

\smallskip

\begin{remarque}
Un modèle est toujours faux : la réalité est trop complexe pour être décrite parfaitement. L'enjeu est que l'approximation apportée par le modèle soit suffisamment fidèle pour répondre à la question posée.  

En pratique, ce sont souvent les spécialistes du domaine (médecine, biologie, économie, etc.) qui proposent ou discutent le modèle.
\end{remarque}

\begin{exemple}
\begin{enumerate}
  \item \textbf{Jeu de pile ou face répété $n$ fois :}
  \[
  E = \{0,1\}^n, \quad 
  \Theta = \,]0,1[\,, \quad
  P_\theta = \big(\mathcal{B}(\theta)\big)^{\otimes n}.
  \]
  Une observation est $X = (X_1,\ldots,X_n) \in \{0,1\}^n$, avec $X_i$ des variables aléatoires indépendantes de loi $\mathcal{B}(\theta)$.

  \item \textbf{Modélisation de la taille d'individus par une loi normale.} On observe sur un échantillon de taille $n$ :
  \[
  E = \mathbb{R}^n, \quad X = (X_1,\ldots,X_n),
  \]
  où les $X_i$ sont des variables aléatoires indépendantes de loi $\mathcal{N}(m,\sigma^2)$, avec
  \[
  \Theta = \{ (m,\sigma^2) : m \in \mathbb{R}, \ \sigma^2 > 0 \} \subset \mathbb{R}^2.
  \]

  \item \textbf{Même population, mais modélisation par une loi à densité sur $[0.5, 2.5]$ :}
  \[
  \Theta = \{ f : [0.5,2.5] \to \mathbb{R}_+ \ \mid f \text{ est une densité} \}.
  \]
  Il s'agit alors d'un \textbf{modèle non paramétrique}.
\end{enumerate}
\end{exemple}

\paragraph{Valeur vraie du paramètre.}
Le point clé de la modélisation est qu'il existe une valeur \emph{vraie} $\theta^* \in \Theta$ telle que :
\[
P_X = P_{\theta^*}.
\]
Mais $\theta^*$ est inconnue.

Dans toute la suite, on notera simplement $\theta$ pour $\theta^*$.

On utilisera les notations :
\[
\mathbb{P}_\theta \quad \text{et} \quad \mathbb{E}_\theta
\]
pour bien marquer la dépendance en $\theta$. On a ainsi :
\[
P(X \in A) = P_X(A) = P_{\theta^*}(A), \qquad P_\theta(X \in A) = P_\theta(A).
\]

\begin{remarque}
On considérera souvent une observation $X$ comme un $n$-échantillon de variables aléatoires i.i.d.\ de loi $Q_\theta$ :
\[
P_\theta = (Q_\theta)^{\otimes n}.
\]
On peut alors définir le modèle statistique comme la famille $\{Q_\theta, \theta \in \Theta\}$.
\end{remarque}

\begin{definition}[Identifiabilité]
Un modèle $\{P_\theta, \theta \in \Theta\}$ est dit \textbf{identifiable} si :
\[
P_{\theta_1} = P_{\theta_2} \quad \Rightarrow \quad \theta_1 = \theta_2.
\]
\end{definition}

\begin{remarque}
L'identifiabilité est une condition \emph{nécessaire} pour pouvoir estimer $\theta$ de manière cohérente : si deux paramètres distincts produisent la même loi, aucune observation ne permet de les distinguer. Dans toute la suite, on supposera le modèle identifiable.

\medskip

\textbf{Exemple.} Considérons $X \sim \mathcal{N}(\mu, \sigma^2)$ avec $\theta = (\mu, \sigma^2)$. La loi de $X$ détermine de façon unique $(\mu, \sigma^2)$ : ce modèle est identifiable. En revanche, si l'on paramétrisait par $(\mu, \sigma)$ et $(\mu, -\sigma)$, les deux valeurs du paramètre donneraient la même loi, et le modèle ne serait plus identifiable.
\end{remarque}

\subsection{Estimateur statistique}

On suppose que le modèle $\{P_\theta, \theta \in \Theta\}$ est identifiable.  
On observe $X$ et on s'intéresse à une quantité d'intérêt $g(\theta)$, où :
\[
g : \Theta \to \mathbb{R}^p, \quad p \geq 1.
\]

\begin{mdframed}
\begin{definition}[Statistique et estimateur]
\begin{itemize}
  \item Une \textbf{statistique} $T(X)$ est une fonction mesurable de $X$ (éventuellement de paramètres connus), qui ne dépend pas de $\theta$.
  \item Un \textbf{estimateur} de $g(\theta)$ est une statistique $\hat{g} = h(X)$ à valeurs dans $\mathbb{R}^p$ (avec $h$ mesurable).
\end{itemize}
\end{definition}
\end{mdframed}

\begin{exemple}
Pour un $n$-échantillon $(X_1,\ldots,X_n)$, les objets suivants sont des statistiques :
\[
\overline{X}_n, \quad X_1, \quad X_2^2 \cdot X_8, \quad 0.
\]
En revanche, $\overline{X}_n + \theta$ \emph{n'est pas} une statistique car elle dépend du paramètre inconnu $\theta$.
\end{exemple}

\begin{remarque}
Si l'on dispose d'une réalisation $x$ de $X$, alors $\hat{g}(x)$ est appelé un \textbf{estimé} de $g(\theta)$.  
Par abus, on note souvent $\hat{g}$ à la fois l'estimateur (variable aléatoire) et son estimé (valeur numérique).
\end{remarque}

\begin{notation}[Estimateur plug-in]
Si $\hat{\theta}$ est un estimateur de $\theta$, un estimateur naturel de $g(\theta)$ est :
\[
g(\hat{\theta}).
\]
On appelle cela un \textbf{estimateur plug-in}.  
Sa qualité dépend de la qualité de $\hat{\theta}$ et de la régularité de $g$.
\end{notation}

\begin{definition}[Risque quadratique]
Le risque quadratique de $\hat{g}$ est défini, pour tout $\theta \in \Theta$, par :
\[
R(\hat{g}, g(\theta)) = \mathbb{E}_\theta \big[ (\hat{g} - g(\theta))^2 \big] \in [0, +\infty].
\]
\end{definition}

\begin{remarque}
\textbf{Lien avec l'inégalité de Markov :}
Par l'inégalité de Markov,
\[
\mathbb{P}_\theta \big( |\hat{g} - g(\theta)| > \varepsilon \big) \leq \frac{R(\hat{g}, g(\theta))}{\varepsilon^2}.
\]
Ainsi, plus le risque quadratique est petit, plus l'estimateur est concentré autour de la valeur vraie.
\end{remarque}

\paragraph{Existe-t-il un estimateur optimal ?}

On pourrait souhaiter qu'il existe un estimateur $\hat{g}$ meilleur que tous les autres pour tout $\theta$, c'est-à-dire :
\[
\exists\, \hat{g} \text{ tel que } \forall\, \theta \in \Theta,\ \forall\, \hat{g}' \text{ estimateur},\quad R(\hat{g}, g(\theta)) \leq R(\hat{g}', g(\theta)).
\]
La réponse est en général \textbf{non}.

\begin{exemple}[Comparaison de deux estimateurs]
Considérons le modèle :
\[
\{ (\mathcal{N}(\theta,1))^{\otimes n}, \ \theta \in \mathbb{R} \}.
\]

Deux estimateurs de $\theta$ :
\[
\hat{\theta}_1 = \overline{X}_n, \qquad \hat{\theta}_2 = 0.
\]

On a :
\[
R(\hat{\theta}_1, \theta) = \frac{1}{n}, \qquad R(\hat{\theta}_2, \theta) = \theta^2.
\]

De plus :
\[
\mathbb{E}_\theta[\hat{\theta}_1] = \theta, \qquad \mathrm{Var}_\theta(\hat{\theta}_1) = \frac{1}{n}.
\]

Conclusion :
\begin{itemize}
  \item $\hat{\theta}_1$ est sans biais avec un risque constant $1/n$.
  \item $\hat{\theta}_2$ est biaisé, mais pour $|\theta| < \frac{1}{\sqrt{n}}$, son risque est inférieur à celui de $\hat{\theta}_1$.
\end{itemize}

Ainsi, même un estimateur « stupide » comme $\hat{\theta}_2$ peut être meilleur que l'estimateur usuel $\overline{X}_n$ pour certaines valeurs de $\theta$. Bien sûr, plus $n$ est grand, plus l'intervalle de $\theta$ où $\hat{\theta}_2$ est meilleur devient petit.
\end{exemple}

\begin{mdframed}
\begin{proposition}[Décomposition biais-variance]
Pour tout $\theta \in \Theta$, on a :
\[
R(\hat{g}, g(\theta)) 
= \big( \mathbb{E}_\theta[\hat{g}] - g(\theta) \big)^2 
+ \mathbb{E}_\theta\!\left[ \big(\hat{g} - \mathbb{E}_\theta[\hat{g}]\big)^2 \right].
\]
Autrement dit :
\[
R(\hat{g}, g(\theta)) = B(\hat{g}, g(\theta))^2 + \mathrm{Var}_\theta(\hat{g}),
\]
où $B(\hat{g}, g(\theta))$ désigne le biais de $\hat{g}$.
\end{proposition}
\end{mdframed}

\begin{definition}[Biais]
Le biais d'un estimateur $\hat{g}$ est défini par :
\[
\theta \mapsto B(\hat{g}, g(\theta)) = \mathbb{E}_\theta[\hat{g}] - g(\theta).
\]
On dit que $\hat{g}$ est \textbf{sans biais} si, pour tout $\theta \in \Theta$ :
\[
\mathbb{E}_\theta[\hat{g}] = g(\theta).
\]
\end{definition}

\begin{proof}
On part de l'identité :
\[
\hat{g} - g(\theta)
= \big( \hat{g} - \mathbb{E}_\theta[\hat{g}] \big)
+ \big( \mathbb{E}_\theta[\hat{g}] - g(\theta) \big).
\]
En élevant au carré :
\[
(\hat{g} - g(\theta))^2
= \big(\hat{g} - \mathbb{E}_\theta[\hat{g}]\big)^2
+ 2\big(\hat{g} - \mathbb{E}_\theta[\hat{g}]\big)\big(\mathbb{E}_\theta[\hat{g}] - g(\theta)\big)
+ \big(\mathbb{E}_\theta[\hat{g}] - g(\theta)\big)^2.
\]
En prenant l'espérance sous $\mathbb{P}_\theta$, le terme croisé disparaît car
\[
\mathbb{E}_\theta\big[\hat{g} - \mathbb{E}_\theta[\hat{g}]\big] = 0.
\]
Ainsi :
\[
\mathbb{E}_\theta\big[(\hat{g} - g(\theta))^2\big]
= \mathbb{E}_\theta\big[(\hat{g} - \mathbb{E}_\theta[\hat{g}])^2\big]
+ \big(\mathbb{E}_\theta[\hat{g}] - g(\theta)\big)^2,
\]
ce qui donne bien $R(\hat{g}, g(\theta)) = \mathrm{Var}_\theta(\hat{g}) + B(\hat{g}, g(\theta))^2$.
\end{proof}

\begin{remarque}
Si $\hat{g}$ est sans biais, alors $B(\hat{g}, g(\theta)) = 0$ et
\[
R(\hat{g}, g(\theta)) = \mathrm{Var}_\theta(\hat{g}).
\]
\end{remarque}

\paragraph{Interprétation :}
\begin{itemize}
  \item Le \textbf{biais} mesure l'erreur systématique faite par l'estimateur (erreur d'approximation).
  \item La \textbf{variance} mesure les fluctuations aléatoires de l'estimateur autour de sa moyenne (erreur d'estimation).
\end{itemize}

\paragraph{Remarques sur les estimateurs sans biais}

Si $\hat{g}$ est sans biais :
\[
R(\hat{g}, g(\theta)) = \mathrm{Var}_\theta(\hat{g}).
\]

On s'intéresse alors aux estimateurs sans biais ayant la variance minimale.  
On appelle un tel estimateur un \textbf{estimateur UVMB} (Uniformément de Variance Minimale parmi les estimateurs Sans Biais).  

Il existe des méthodes (exhaustivité, complétude, théorème de Lehmann-Scheffé) pour caractériser directement de tels estimateurs.

\begin{remarque}
Un estimateur sans biais n'est pas toujours le meilleur choix :  
un estimateur \textbf{avec biais} peut avoir un risque quadratique plus petit.  
De plus, il n'existe pas toujours d'estimateur sans biais pour une quantité donnée.
\end{remarque}

\begin{exemple}[Impossibilité d'un estimateur sans biais]
Soit $X \sim \mathcal{E}(\theta)$ une loi exponentielle de paramètre $\theta > 0$.  
Supposons qu'il existe un estimateur sans biais $\hat{\theta} = h(X)$ de $\theta$.  
Alors, pour tout $\theta > 0$, il faudrait :
\[
\int_0^\infty h(x)\, \theta e^{-\theta x}\, dx = \theta.
\]
Or, cette condition n'admet pas de solution mesurable $h$, donc un tel estimateur n'existe pas.
\end{exemple}

\begin{remarque}
Même si $\hat{\theta}$ est un estimateur sans biais de $\theta$, en général $g(\hat{\theta})$ n'est \emph{pas} un estimateur sans biais de $g(\theta)$ (sauf si $g$ est affine).  
Ainsi, la propriété « sans biais » n'est pas stable par transformation non linéaire.
\end{remarque}





\newpage

\section{Propriétés asymptotiques}

On se place dans le cadre de $n$ observations.  
On note $\hat{g}_n$ un estimateur de $g(\theta)$.

\smallskip

\begin{mdframed}
\begin{definition}[Consistance]
On dit que la suite $(\hat{g}_n)_{n \in \mathbb{N}}$ est \textbf{consistante} si, pour tout $\theta \in \Theta$,
\[
\hat{g}_n \xrightarrow[\mathbb{P}_\theta]{n \to +\infty} g(\theta).
\]
Elle est dite \textbf{fortement consistante} si, pour tout $\theta \in \Theta$,
\[
\hat{g}_n \xrightarrow[\text{p.s.}]{n \to +\infty} g(\theta).
\]
\end{definition}
\end{mdframed}

\begin{remarque}
\begin{itemize}
    \item La convergence doit avoir lieu pour \emph{tout} $\theta \in \Theta$.
    \item La consistance est une propriété « faible ». S'il existe un estimateur consistant, il en existe une infinité : si $\hat{g}_n$ est consistant et $(a_n)$ vérifie $a_n \to 1$, alors $a_n \hat{g}_n$ est aussi consistant.
    \item Les estimateurs non consistants sont à exclure : ils ne convergent pas vers la vraie valeur, quelle que soit la taille de l'échantillon.
\end{itemize}
\end{remarque}

\begin{exemple}[Consistance de la moyenne empirique]
Soit $(X_1, \ldots, X_n)$ un $n$-échantillon i.i.d.\ de loi $P_\theta$, avec $\mathbb{E}_\theta[X_1] = \mu(\theta) < +\infty$.  
La loi des grands nombres assure que
\[
\overline{X}_n = \frac{1}{n}\sum_{i=1}^n X_i \xrightarrow[\text{p.s.}]{n \to +\infty} \mu(\theta).
\]
Ainsi, $\overline{X}_n$ est un estimateur fortement consistant de $\mu(\theta)$.
\end{exemple}

\begin{definition}[Estimateur asymptotiquement sans biais]
On suppose que pour tout $\theta \in \Theta$ et tout $n \in \mathbb{N}$,
\[
\mathbb{E}_\theta\!\left[\|\hat{g}_n\|\right] < +\infty.
\]
On dit que $\hat{g}_n$ est \textbf{asymptotiquement sans biais} si, pour tout $\theta \in \Theta$,
\[
\mathbb{E}_\theta[\hat{g}_n] \xrightarrow{n \to +\infty} g(\theta).
\]
\end{definition}

\begin{remarque}
Asymptotiquement sans biais n'implique pas consistant, et réciproquement. En revanche, un estimateur sans biais et consistant est bien asymptotiquement sans biais.
\end{remarque}

\smallskip

\begin{mdframed}
\begin{definition}[Normalité asymptotique]
On dit que $\hat{g}_n$ est \textbf{asymptotiquement normal} si, pour tout $\theta \in \Theta$, il existe $\Sigma(\theta) > 0$ tel que
\[
\sqrt{n}\left(\hat{g}_n - g(\theta)\right) \xrightarrow[\text{loi}]{n \to +\infty} \mathcal{N}(0, \Sigma(\theta)).
\]
La quantité $\Sigma(\theta)$ est appelée la \textbf{variance asymptotique} de $\hat{g}_n$.
\end{definition}
\end{mdframed}

\begin{remarque}
\begin{itemize}
    \item Si $\hat{g}_n$ est asymptotiquement normal, alors $\hat{g}_n$ est consistant.
    \item La normalité asymptotique est fondamentale en pratique : elle permet de construire des intervalles de confiance et des tests asymptotiques.
    \item Plus généralement, on peut remplacer $\sqrt{n}$ par $a_n$ avec $a_n \to +\infty$. On dit alors que $1/a_n$ est la \textbf{vitesse de convergence} asymptotique.
\end{itemize}
\end{remarque}

\smallskip

\begin{mdframed}
\begin{proposition}
Soit $\hat{g}_n$ un estimateur asymptotiquement normal de $g(\theta)$ de variance asymptotique $\Sigma(\theta)$.
Soit $\hat{\Sigma}_n$ un estimateur consistant de $\Sigma(\theta)$, c'est-à-dire :
\[
\hat{\Sigma}_n \xrightarrow[\mathbb{P}_\theta]{n \to +\infty} \Sigma(\theta) \quad \forall\, \theta \in \Theta.
\]
Alors
\[
\frac{\sqrt{n}\left(\hat{g}_n - g(\theta)\right)}{\sqrt{\hat{\Sigma}_n}} \xrightarrow[\text{loi}]{n \to +\infty} \mathcal{N}(0, 1).
\]
\end{proposition}
\end{mdframed}

\begin{remarque}
Ce résultat est particulièrement utile en pratique : on remplace la variance asymptotique inconnue $\Sigma(\theta)$ par un estimateur $\hat{\Sigma}_n$, et la convergence en loi est préservée. Cela permet de construire des intervalles de confiance asymptotiques sans connaître $\theta$.
\end{remarque}

\begin{proof}
On écrit
\[
\frac{\sqrt{n}\left(\hat{g}_n - g(\theta)\right)}{\sqrt{\hat{\Sigma}_n}}
= \frac{\sqrt{n}\left(\hat{g}_n - g(\theta)\right)}{\sqrt{\Sigma(\theta)}}
\cdot \frac{\sqrt{\Sigma(\theta)}}{\sqrt{\hat{\Sigma}_n}},
\]
avec :
\begin{itemize}
    \item $\displaystyle\frac{\sqrt{n}\left(\hat{g}_n - g(\theta)\right)}{\sqrt{\Sigma(\theta)}} \xrightarrow[\text{loi}]{n \to +\infty} \mathcal{N}(0,1)$ par hypothèse,
    \item $\displaystyle\frac{\sqrt{\Sigma(\theta)}}{\sqrt{\hat{\Sigma}_n}} \xrightarrow[\mathbb{P}_\theta]{n \to +\infty} 1$ par consistance de $\hat{\Sigma}_n$.
\end{itemize}
Le \textbf{lemme de Slutsky} permet de conclure : le produit d'une suite convergeant en loi vers $Z$ et d'une suite convergeant en probabilité vers $1$ converge en loi vers $Z$.
\end{proof}

\section{Estimation de $\theta$ à partir d'un estimateur de $g(\theta)$}

\textbf{Question :} On dispose de $\hat{g}_n$, estimateur de $g(\theta)$. Comment en déduire un estimateur de $\theta$ ? Quelles propriétés possède-t-il ?

\smallskip

L'outil fondamental pour répondre à cette question est la \textbf{méthode delta}, que l'on énonce d'abord dans sa forme directe.

\begin{mdframed}
\begin{proposition}[Méthode delta]
Soit $\Theta$ un ouvert de $\mathbb{R}$. Soit $\hat{g}_n$ un estimateur de $g(\theta)$ tel que pour tout $\theta \in \Theta$,
\[
\sqrt{n}\left(\hat{g}_n - g(\theta)\right) \xrightarrow[\text{loi}]{n \to +\infty} \mathcal{N}(0, \sigma^2).
\]
Soit $\phi : g(\Theta) \to \mathbb{R}$ une fonction de classe $\mathcal{C}^1$. Alors, pour tout $\theta \in \Theta$,
\[
\sqrt{n}\left(\phi(\hat{g}_n) - \phi(g(\theta))\right) \xrightarrow[\text{loi}]{n \to +\infty} \mathcal{N}\!\left(0,\, \sigma^2 \left(\phi'(g(\theta))\right)^2\right).
\]
\end{proposition}
\end{mdframed}

\begin{remarque}
La méthode delta permet de propager la normalité asymptotique au travers d'une transformation régulière. Elle est très utilisée pour construire des estimateurs et quantifier leur précision.
\end{remarque}

\smallskip

On applique maintenant la méthode delta à la fonction $\phi = g^{-1}$, ce qui donne la version suivante.

\begin{mdframed}
\begin{proposition}[Méthode delta inverse]
Soit $\Theta$ un intervalle ouvert de $\mathbb{R}$.
\begin{itemize}
    \item Soit $\hat{g}_n$ estimateur de $g(\theta)$ tel que pour tout $\theta \in \Theta$,
    \[
    \sqrt{n}\left(\hat{g}_n - g(\theta)\right) \xrightarrow[\text{loi}]{n \to +\infty} \mathcal{N}(0, \sigma^2).
    \]
    \item On suppose que $g : \Theta \to g(\Theta)$ est un $\mathcal{C}^1$-difféomorphisme (bijective, $\mathcal{C}^1$, et $g^{-1} \in \mathcal{C}^1$). On définit :
    \[
    \hat{\theta}_n = g^{-1}(\hat{g}_n).
    \]
\end{itemize}
Alors :
\begin{itemize}
    \item $\hat{\theta}_n$ est consistant pour l'estimation de $\theta$.
    \item Pour tout $\theta \in \Theta$,
    \[
    \sqrt{n}\left(\hat{\theta}_n - \theta\right) \xrightarrow[\text{loi}]{n \to +\infty} \mathcal{N}\!\left(0,\, \frac{\sigma^2}{(g'(\theta))^2}\right).
    \]
\end{itemize}
\end{proposition}
\end{mdframed}

\begin{proof}
\begin{itemize}
    \item \textbf{Bonne définition de $\hat{\theta}_n$ :} $g$ est continue et bijective de $\Theta$ sur $g(\Theta)$, qui est un intervalle ouvert. Comme $\hat{g}_n \xrightarrow{\mathbb{P}_\theta} g(\theta) \in g(\Theta)$, on a
    \[
    \mathbb{P}\!\left(\hat{g}_n \in g(\Theta)\right) \xrightarrow{n \to +\infty} 1.
    \]
    Donc $\hat{\theta}_n = g^{-1}(\hat{g}_n)$ est bien défini avec une probabilité tendant vers $1$.

    \item \textbf{Consistance de $\hat{\theta}_n$ :} $g^{-1}$ est continue et $\hat{g}_n \xrightarrow{\mathbb{P}_\theta} g(\theta)$, donc par composition,
    \[
    \hat{\theta}_n = g^{-1}(\hat{g}_n) \xrightarrow{\mathbb{P}_\theta} g^{-1}(g(\theta)) = \theta.
    \]

    \item \textbf{Normalité asymptotique :} on applique la méthode delta à $\phi = g^{-1}$. On rappelle que
    \[
    (g^{-1})'(y) = \frac{1}{g'(g^{-1}(y))},
    \]
    d'où $(g^{-1})'(g(\theta)) = \dfrac{1}{g'(\theta)}$.

    Par la méthode delta :
    \[
    \sqrt{n}\left(\hat{\theta}_n - \theta\right) = \sqrt{n}\left(g^{-1}(\hat{g}_n) - g^{-1}(g(\theta))\right) \xrightarrow[\text{loi}]{n \to +\infty} \mathcal{N}\!\left(0,\, \frac{\sigma^2}{(g'(\theta))^2}\right).
    \]
\end{itemize}
\end{proof}
